\documentclass[a4paper,10pt]{article}
\usepackage{graphicx}

\usepackage[utf8]{inputenc}
\usepackage[portuguese,brazil]{babel}
\usepackage[T1]{fontenc}

\usepackage{indentfirst}
\usepackage{url}


\begin{document}


\begin{titlepage}

\begin{minipage}{0.2\linewidth}
 \includegraphics[]{./minerva.png}
\end{minipage}
\begin{minipage}{0.8\linewidth}
 \textbf{Universidade Federal do Rio de Janeiro}\\
 Instituto de Matemática\\
 Departamento de Ciência da Computação\\
 \rule{0.8\linewidth}{0.5mm}\\
 Rio de Janeiro, RJ - Brasil
\end{minipage}

\begin{center}

\vspace{2cm}

\Large
Trabalho de Simulaçao: Implementação e análise de um simulador.

\vspace{1cm}

\large

\_\_\_\_\_\_\_\_\_\_\_\_\_\_\_\_\_\_\_\_\_\_\_\_\_\_\_\_\_\_\_\_\_\_\_\_\_\_\_\\
Bruno C. Buss\\(Implementação, documentação do código e relatório final)\\

\vspace{0.5cm}

\_\_\_\_\_\_\_\_\_\_\_\_\_\_\_\_\_\_\_\_\_\_\_\_\_\_\_\_\_\_\_\_\_\_\_\_\_\_\_\\
Felipe P. Martinez\\(Implementação, documentação do código e relatório final)\\

\vspace{0.5cm}

\_\_\_\_\_\_\_\_\_\_\_\_\_\_\_\_\_\_\_\_\_\_\_\_\_\_\_\_\_\_\_\_\_\_\_\_\_\_\_\\
Rafael O. Lopes\\(Implementação, documentação do código e relatório final)\\

\vspace{0.5cm}

\_\_\_\_\_\_\_\_\_\_\_\_\_\_\_\_\_\_\_\_\_\_\_\_\_\_\_\_\_\_\_\_\_\_\_\_\_\_\_\\
Yanko G. Oliveira\\(Implementação, documentação do código e relatório final)\\

\vspace{0.5cm}

\vspace{1cm}

Relatório gerado em \today

\normalsize
\end{center}

\vfill

\begin{flushright}
Disciplina: Avaliação e Desempenho 2010/2\\
Professor: Paulo Henrique de Aguiar Rodrigues\\
\end{flushright}

\vspace{2cm}

\end{titlepage}

\pagebreak

\tableofcontents
\pagebreak

\section{Introdução}
\subsection{Funcionamento Geral do Simulador}
    O simulador é executado por linha de comando (comando padrão: "./cmulator"). Este aguarda diversos parâmetros que podem ser passados pela linha de comando ou, caso isso não seja feito, são pedidos em tempo de execução. Tais parâmetros são:

\begin {itemize}
\item \textbf{Modo:} O usuário escolhe qual o parâmetro de execução da simulação, se será replicativo ou batch. Na linha de comando: \emph{--modo} ou \emph{-m}. Opções possíveis: \emph{"Batch"} ou \emph{"Replicativo"};

\item \textbf{Quantidade de rodadas:} O usuário decide a quantidade de rodadas; esta deve ser superior a 10, pois apenas assim temos um valor assintótico na tabela t-student para o cálculo do intervalo de confiança. Na linha de comando: \emph{--n\_rodadas} ou \emph{-n}. Opções possíveis: números inteiros;

\item \textbf{Tamanho da rodada:} Determina a quantidade de fregueses típicos em cada rodada. Na linha de comando: \emph{--t\_rodada} ou \emph{-r}. Opções possíveis: números inteiros;

\item \textbf{Tamanho da fase transiente:} Determina a quantidade de fregueses típicos a serem considerados pertencentes à fase transiente, de forma a considerarmos apenas dados coletados após o sistema estar em equilíbrio. Na linha de comando: \emph{--t\_transiente} ou \emph{-t}. Opções possíveis: números inteiros;

\item \textbf{Tipo da fila 1:} O tipo da fila pode ser FCFS(fila), ou LCFS(pilha). Na linha de comando: \emph{--fila\_1} ou \emph{-1}. Opções possíveis: \emph{"FCFS"} ou \emph{"LCFS"};

\item \textbf{Tipo da fila 2:} Assim como a fila 1, a fila 2 pode ser FCFS ou LCFS. Na linha de comando: \emph{--fila\_2} ou \emph{-2}. Opções possíveis: \emph{"FCFS"} ou \emph{"LCFS"};

\item \textbf{Taxa $\lambda$:} Determina a taxa de chegada de chegada de fregueses no sistema. Na linha de comando: \emph{--tx\_lambda} ou \emph{-l}. Opções possíveis: números reais;

\item \textbf{Taxa $\mu$:} Determina a taxa de serviço de fregueses no sistema. Na linha de comando: \emph{--tx\_mi} ou \emph{-m}. Opções possíveis: números reais.
\end {itemize}

    Além desses parâmetros, que são obrigatoriamente pedidos, existem outros parâmetros opcionais para o simulador, a serem passados pela linha de comando:

\begin {itemize}
\item \textbf{Semente para gerador de chegadas:} semente utilizada para o gerador de números pseudo-aleatórios responsável por gerar os tempos de chegada. Na linha de comando: \emph{--seed\_gerador\_chegadas} ou \emph{-c}. Opções possíveis: números reais;

\item \textbf{Semente para gerador de tempos de serviço:} semente utilizada para o gerador de números pseudo-aleatórios responsável por gerar os tempos de serviço. Na linha de comando: \emph{--seed\_gerador\_tempo\_servico} ou \emph{-x}. Opções possíveis: números reais;

\item \textbf{Sobre:} imprime informações sobre o simulador e seus autores. Na linha de comando: \emph{--sobre} ou \emph{-s}. Não recebe parâmetros;

\item \textbf{Ajuda:} imprime todos os comandos possíveis e suas descrições. Na linha de comando: \emph{--ajuda} ou \emph{-a}. Não recebe parâmetros;

\item \textbf{Modo verborrágico:} roda o simulador imprimindo diferentes níveis de detalhe sobre as operações sendo executadas internamente. Na linha de comando: \emph{--verbose} ou \emph{-v} Opções: 0,1 ou 2;
\end {itemize}

    Após a entrada de parâmetros, o simulador se inicia. Este utiliza eventos discretos, mas tempo contínuo, logo, os eventos acontecem em instantes de tempo com valores numéricos reais. A cada rodada, o simulador fica em loop aguardando eventos acontecerem. O loop se encerra quando a quantidade de fregueses determinada pelo parâmetro \emph{"tamanho da rodada"} for servida. Os eventos que acontecem são:

\begin {itemize}
\item \textbf{Chegada de um freguês na fila 1:} Um freguês chega e é inserido na fila 1, de acordo com a disciplina de atendimento escolhida - caso seja FCFS, o mesmo é inserido no final e, caso seja LCFS, é adicionado no início.

\item \textbf{Término do serviço do freguês no servidor:} No momento em que um freguês termina seu atendimento no servidor, é verificado qual sua fila original. Caso seja a fila 1,
ele é movido para a fila 2 (sempre respeitando a disciplina de atendimento). Se o freguês veio da fila 2, dados estatísticos são coletados e o mesmo é retirado do sistema.
\end {itemize}

    Após a tratar o evento ocorrido, verifica-se se o servidor está vazio. Caso esteja, é inserido no servidor um freguês, de acordo com as regras determinadas; neste caso, dando prioridade aos fregueses da fila 1 e sendo os da fila 2 servidos apenas quando a primeira estiver vazia.

    Após a execução de todas as rodadas, os dados estatísticos coletados são analisados e é calculado o intervalo de confiança para cada variável aleatória. O programa se encerra imprimindo na tela a linha de comando com todos os parâmetros utilizados (mesmo que estes tenham sido entrados em tempo de execução), incluindo as sementes geradoras para, caso seja desejado, executar mais uma vez a exata mesma rotina.

\subsection{Linguagem de Programação Utilizada}

    Foi utilizada a linguagem de programação C++. O compilador utilizado para desenvolver o sistema foi o GCC 4.4, presente em qualquer distribuição atualizada do Linux, de forma que o sistema pode ser compilado em qualquer ambiente que possua o GCC 4.4 ou superior instalado.

\emph{Nota: o GCC padrão incluso no pacote MingW para compilação no ambiente Windows não possui algumas funções utilizadas pelo simulador (neste caso, as relacionadas à geração de números pseudo-aleatórios). Sendo assim, para que o mesmo seja compilado nesta plataforma, é necessário o uso de GCCs modificados e/ou ferramentas como o Cygwin.}

    O principal fator para escolha da linguagem foi a possibilidade de lidar com orientação a objetos, tornando o código mais sucinto e de melhor entendimento. Um outro fator é o fato de que C++ é a linguagem com que os membros do grupo estão melhor familiarizados, proporcionando uma maior facilidade na hora da implementação.

    O programa foi completamente desenvolvido e testado em ambiente Linux, não havendo garantias de seu funcionamento em ambiente Windows.

\subsection{Estruturas Internas Utilizadas}

    Todo o código do simulador está comentado e foi documentado utilizando a ferramenta \emph{DoxyGen}, estando a documentação gerada anexa. Nesta seção, faremos uma breve análise da implementação e estruturas utilizadas.

    Todas as estruturas do simulador estão no namespace \emph{TrabalhoAD} permitindo não só maior organização do código do próprio trabalho mas também caso eventualmente este seja tilizado juntamente a outras aplicações. Nele há as classes \emph{Evento, DistExponencial, Fregues} e \emph{Simulador}, as structs \emph{AmostragemFila} e \emph{ResultadosConsolidados}, além dos operadores sobrecarregados \emph{$>$} e \emph{$>=$} e das enumerações \emph{ETipo} e \emph{TipoFila}.

    Os operadores sobrecarregados são responsáveis por comparar dois eventos do simulador, sendo o maior aquele com o tempo em que ocorreu de valor mais alto. Os enums \emph{ETipo} e \emph{TipoFila} listam, respectivamente, os tipos de evento possíveis (se um novo freguês chegou ao sistema, ou se seu serviço terminou de ser executado) e o regime de serviço das filas (FCFS ou LCFS).

    Descrição breve das classes:

\begin {itemize}
\item \textbf{class Evento:} representa um evento ocorrido no sistema, descrevendo seu tipo o tempo em que ocorrerá;

\item \textbf{class DistExponencial:} responsável pelo gerador pseudo-aleatório de tempos entre chegadas exponencialmente distribuídas. O método de geração é através da função inversa a PDF da distribuição exponencial. Além disso, a semente inicial é guardada, para que a simulação possa ser repetida de forma precisa;

\item \textbf{class Fregues:} representa um freguês que chegou durante a simulação. É a responsável por gravar todos os dados da simulação (o tempo em que entrou e saiu das filas, o tempo de espera, qual fila está atualmente, a qual rodada pertence etc);

\item \textbf{class Simulador:} responsável por executar a simulação propriamente dita, coletar e armazenar os resultados de cada rodada e apresentar o resultado final da simulação.
\end {itemize}

    Além destas, temos os structs:

\begin {itemize}
\item \textbf{struct AmostragemFila:} unifica todos os dados de uma determinada fila. Nela constam as quantidades de pessoas na fila, a soma dos tempo de espera, no servidor etc e seus quadrados, para posterior análise estatística;

\item \textbf{struct ResultadosConsolidados:} reúne todos os dados coletados das \emph{duas} filas (através de duas \emph{AmostragemFila}) e da quantidade de amostras analisadas.
\end {itemize}

    Um detalhe importante sobre a implementação do simulador é a utilização de geradores de números pseudo-aleatórios diferentes do padrão (como citado na seção "Linguagem de Programação Utilizada"): \emph{drand48\_r\(\)} e \emph{srand48\_r\(\)}. Esta escolha foi feita pois além de retornarem um número pseudo-aleatório uniformemente distribuído entre [0.0, 1.0), estas funções são re-entrantes, o que evita problemas nos casos em que temos mais de um classe deste tipo gerando tempos de chegada concomitantemente.


% Aqui comentamos sobre as classes do projeto, e alguns detalhes, como a utilização do rand do linux, ou de
\subsection{Implementação da Lista de Eventos}
% Creio que essa seção pode morrer, pois a lista de eventos já foi mencionada antes. Talvez possamos desenvolver melhor a parte de eventos lá.
\subsection{Geração das Variáveis Aleatórias}
% Aqui dizemos as variáveis aleatórias que temos, e como elas são calculadas(mencionar a inversa da exponencial etc)
\subsection{Coleta de Estatísticas}
% Falamos sobre a coleta de estatísticas e cálculo delas(Variância e IC).
\subsection{Método de análise}

    O simulador foi construído para funcionar de maneira genérica, sendo assim, a única parte fixa é o formato das filas do sistema (fila 1 prioritária sobre fila 2, ambas sendo servidas por um mesmo servidor), enquanto as disciplinas de serviço podem ser escolhidas pelo usuário. Com isso, também foi planejado para aceitar ambos os métodos de análise: Batch e Replicativo.

    No método Batch, diversas rodadas são executadas a partir de uma mesma semente. Desta forma, precisamos estimar a fase transiente apenas uma vez, utilizando apenas um sistema de \emph{"coloração"} para definir quais dados devem ser coletados. Os clientes gerados em uma determinada rodada possuem um identificador da mesma e, quando saem do sistema, apenas são considerados para as estatísticas aqueles pertencentes à rodada atual de simulação. Caso a rodada acabe e ainda haja fregueses na fila, estes permanecem fazendo parte da simulação, mas seus dados não são contados na rodada seguinte.

    No método Replicativo, há uma semente gerada para cada rodada. Sendo assim, há também uma fase transiente para ser estimada. Assim como em uma rodada anterior do método Batch, os fregueses não atendidos que estavam na fila de espera em uma rodada continuam na fila, mas não são computados nas estatísticas da rodada referente.


% Aqui mencionamos que decidimos testar os dois métodos: Batch e Replicativo. E também falar um pouco sobre eles.
% Acredito que essa seção pode morrer.
%\subsection{Parâmetros Utilizados}
% Essa sub-seção eu peguei do relatório do Jonas. No projeto deles, eles determinaram um valor para alguns parâmetros que nós tomamos como fixos. Creio que isso pode morrer.
\subsection{Tempo Gasto Para Simulação}
% Tempo que gastamos nos testes

\pagebreak

\section{Teste de Correção}
% Aqui temos que provar que está correto.
% Para fazer isso, devemos comparar os resultados da execução determinística com os resultados da execução exponencial.
\pagebreak

\section{Estimativa da fase transiente}
% Benchmark do Buss
\pagebreak

\section{Tabelas com resultados e comentários pertinentes}
% Tabelas, tabelas e tabelas. =)
\pagebreak

\section{Otimização}
% O que falar exatamente aqui?
\pagebreak

\section{Conclusão}

\pagebreak

\end{document}
